% !TEX root = ../notes_sovereignai.tex
% Chapter 1: The Rig, a build manual.
% One microsection per part, in assembly order: name the part, say what its
% datasheet numbers buy the rig, then fit it. Rationale, alternatives, and
% background live in marginnotes. Every part cites its datasheet. No prices in the
% part sections: the only figure is the day one total in the chapter opening, so
% the chapter cannot go stale line by line.
%
% The parts that are populated onto the board while it lies flat, the processor
% and cooler, the memory, and the drives, are SUBSECTIONS of The Motherboard,
% because that is one bench job in the assembly order rather than four. Its "fit
% it" paragraph, "Work on it flat", is the lead in to them. Parts that attach to
% the finished board or hang off the frame stay top level sections.
%
% Scope stops at the wired machine. Firmware setup and the enumeration probe are
% the first thing software does to the rig, so they live in the Engine chapter,
% which owns everything from the first power-on onward. Nothing here switches
% the machine on.
%
% The main column builds a ONE CARD machine, start to finish. Everything about
% growing to four cards lives in the closing "Cards Two to Four" section, which is
% the potential next extension, not part of today's build. Do not let expansion
% mechanics back into the part sections.
%
% Two closing sections follow the build, both alternative paths rather than work
% the reader does today: "Cards Two to Four" scales the rig up, "The Minimal Viable
% Rig" scales it down to a Raspberry Pi for a reader who wants every step of this
% book without the silicon. The Pi section keeps the same three newthought shape
% and the same scope rule: it stops at the assembled, unpowered board.
%%%%%%%%%%%%%%%%%%%%%%%%%%%%%%%%%%%%%%%%%%%%%%%%%%%%%%%%%%%%%%%%%%%%%%

\index{rig}\index{GPU}\index{graphics card|see{GPU}}\index{RTX 3090}\index{PCIe}%
\index{power supply}\index{sovereign AI}\index{VRAM}\index{video memory|see{VRAM}}
\chapter{The Rig}
\printchapterversion{01_hardware}
\newthought{Owning the Metal.}
\vspace{1.5em}

\newthought{This chapter builds one machine}: An open frame rig with a single
RTX~3090 and its 24~GB, which is enough to serve a capable model for private usage.
The rest of the setup is designed in a way so it is able to grow with your needs.
Order the parts in Table~\ref{tab:bom}\index{bill of materials} and fit them in the order of the sections
that follow, until the last cable is seated and the machine is ready to be switched
on. Day one is roughly 2500~euro and a few hours plugging.

\newthought{The parts list} is the Ryzen configuration of a proven quad~3090 local
AI home server, priced at 3000~EUR when it was built by me in 2026.

% Full width table*: the tabular spans the text block AND the margin column, but the
% caption still sets in the margin, below the table. Nothing else may occupy the
% margin at that height: a marginnote or a citation footnote landing there prints
% straight through the caption with no LaTeX warning. That is why the cited
% paragraph sits ABOVE the table, so its footnote clears the caption.
% Column widths are fractions of \linewidth, which is the full measure in here.
\vspace{2\baselineskip}
\begin{table*}[!ht]
\centering
\begin{tabular}{p{0.24\linewidth}p{0.58\linewidth}p{0.10\linewidth}}
    \toprule
    \textbf{Component} & \textbf{Part} & \textbf{Pieces} \\
    \midrule
    Frame & open GPU rack frame & 1 \\
    Motherboard & Gigabyte B650 Eagle AX & 1 \\
    CPU & AMD Ryzen 5 9600X, AM5 & 1 \\
    CPU cooler & AM5 tower air cooler & 1 \\
    Memory & G.SKILL Ripjaws S5 64\,GB DDR5-6000 & 1 \\
    Storage & NVMe Gen4 x4 2\,TB, M.2 2280 & 1 \\
    GPU & RTX~3090 24\,GB, used & 1 (4) \\
    Risers, x1 & PCIe 3.0 x1 riser, cards 2 to 4 & 0 (3) \\
    Power supply & ASRock PG-1600G, 1600\,W & 1 \\
    \bottomrule
\end{tabular}
\vspace{1.5em}
\caption{The bill of materials, in the order the parts go together. A bracketed
number is what this board, frame, and supply will carry later.} \label{tab:bom}
\end{table*}

% ==============================================================================================
% THE MOTHERBOARD
% ==============================================================================================

\newpage
\section{The Motherboard}\index{motherboard}\index{mainboard|see{motherboard}}\index{PCIe!lanes}

\newthought{A Gigabyte B650 Eagle AX}\index{Gigabyte B650 Eagle AX}, a mainstream board that support our design
decisions \cite{gigabyteB650EagleAX}:

\marginnote{\newthought{A slow slot} therefore costs seconds per model swap, while
this does not touch inference, we still care.}

\begin{itemize}[topsep=0.25em, itemsep=0.4em, parsep=0em]
    \item To support the processor we will be choosing, the socket must be AM5\index{AM5 socket}.
    \item For the GPU the top x16 slot\index{PCIe!x16 slot} is best when running the sixteen PCIe lanes,
    short for peripheral component interconnect express, straight from the processor
    rather than through the chipset\index{chipset}. A model's whole weight file crosses that link
    every time it loads, so model swapping\index{model swapping} will benefit from this wiring.
    \item Four slots have to be x16 sized, so we can grow the server in terms of GPU
    later on. Total VRAM\index{VRAM}, the video memory on the cards themselves, is the ceiling
    on model size, and we can extend that by up to 4 GPUs with this characteristics.
    \item For storage we need one M.2 socket\index{M.2 socket} wired PCIe~4.0~x4 or better straight
    from the processor, and this board has three, so a second drive later costs
    nothing today. That one drive carries everything the server accumulates, model
    weights, container volumes, embeddings, documents, and everything else data.
\end{itemize}

\begin{figure}[!ht]
\centering
\includegraphics[width=0.45\linewidth]{01_board_bare_flat_hl}
\caption{Gigabyte B650 Eagle AX}\label{fig:board}
\end{figure}

\newthought{Work on it flat.} Lay the board on its own box, on the antistatic bag\index{antistatic bag} it
shipped in, and populate it there: Processor, cooler, memory, and the drive all go
in. This manual guides you step by step.

\marginnote{\newthought{Read the manuals.} Every part ships one, and it holds the
pin maps, socket order, and slot wiring that no build guide reproduces correctly.}

% ==============================================================================================
% THE PROCESSOR AND COOLER
% ==============================================================================================

\vspace{1.5em}
\needspace{9\baselineskip}
\subsection{The Processor and Cooler}

\newthought{An AMD Ryzen 5 9600X}\index{Ryzen 5 9600X}\index{CPU}, paired with an AM5 tower air cooler\index{air cooler}, a modest chip
that supports our design decisions\cite{amdRyzen9600X}.

\marginnote{\newthought{Why a tower and not a low profile block.} An open frame has
no lid, so the height a tower needs costs nothing here, and there is no case airflow
for a flat block to borrow either. And towers are also more efficient 
and quieter than flat blocks.}

\begin{itemize}[topsep=0.25em, itemsep=0.4em, parsep=0em]
    \item For inference we need very little compute on the CPU. Six Zen~5\index{Zen 5} cores, twelve threads, 3.9~GHz base and up to 5.4 boost, 32~MB
    of L3, and a 65~watt default thermal design power\index{thermal design power}, is perfect for this.  The GPUs carry the real work while the processor only feeds them and runs the container
    workloads layered on top.
    \item For the top slot to run a real x16, PCIe has to come from the processor
    itself rather than through the chipset.
    \item For the display we want integrated graphics\index{integrated graphics} on the chip itself, so the
    monitor never has to occupy a slot, and a 3090 stays free for the inference and AI model load.
    \item For the cooler we need an AM5 tower air cooler. The fan on the front of the fin stack sits over
    the first memory slot, so a fan with adjustable positioning is smart.
\end{itemize}

\begin{figure}[!ht]
\centering
\includegraphics[width=0.45\linewidth]{02_board_cpu_seated_hl}
\caption{AMD Ryzen 5 9600X, AM5}\label{fig:cpu}
\end{figure}

\begin{figure}[!ht]
\centering
\includegraphics[width=0.45\linewidth]{05_cooler_top_hl}
\caption{AM5 tower air cooler}\label{fig:cooler}
\end{figure}


\newthought{Seat it} by matching the small triangle on the chip's corner to the
triangle on the socket, lower it flat without pressure, close the retention arm
(Figure~\ref{fig:cpu}).




\newthought{Thermal paste}\index{thermal paste} is often pre-applied to the cooler, but if it is not, apply a small pea sized drop to the
    center of the chip, and let the cooler spread it when you mount it.
    Too much paste behaves like a thermal insulator, not a conductor, so don't overdo it at this point.

\newthought{Where the fan plugs in.} The cooler's fan goes on the four
pin header marked \texttt{CPU\_FAN}, right beside the socket. It is the one header
the firmware watches, and a fan anywhere else leaves it reading zero rpm on the
processor, which most boards refuse to boot on.

\marginnote{\newthought{The five fan headers.}\index{fan header} Each is four pin and good for
24~watt.
\begin{itemize}[leftmargin=*, topsep=0.25em, itemsep=0.1em, parsep=0em]
    \item \texttt{CPU\_FAN}: the cooler's fan goes here.
    \item \texttt{CPU\_OPT}: a second CPU fan or pump, right next to it, useful if
    your tower cooler is a dual fan model.
    \item \texttt{SYS\_FAN1}, \texttt{SYS\_FAN2}, \texttt{SYS\_FAN3\_PUMP}: the
    three for the fans you zip tie to the frame rail later.
\end{itemize}}





% ==============================================================================================
% MEMORY
% ==============================================================================================

% No \newpage here: the board section runs long enough that the processor
% subsection breaks on its own, and memory flows onto the page that break opens.
% Memory has to fit that page whole: if its last paragraph spills, storage's
% \newpage below leaves the spill alone on a page of its own. That is why this
% subsection carries no \vspace of its own, and why its prose is cut this tight.
\needspace{9\baselineskip}
\subsection{Memory or RAM}

\newthought{A G.SKILL Ripjaws S5 kit, 64~GB in two 32~GB modules}\index{G.SKILL Ripjaws S5}\index{memory}\index{RAM|see{memory}}, part number
F5-6000J3636F32GX2-RS5K\cite{gskillRipjawsS5}:

\begin{itemize}[topsep=0.25em, itemsep=0.4em, parsep=0em]
    \item RAM kees a weight file warm between loads, plus room
    for the containers layered on later.
    \item For speed the kit runs DDR5-6000\index{memory!DDR5}, double data rate, at CL36-36-36-96\index{memory!CAS latency}
    column address strobe latency and 1.35~volt. Its AMD EXPO profile\index{EXPO profile}, extended
    profiles for overclocking, is where we will take this from in the configuration later on.
    \item We stay at two modules as filling all four slots pushes the memory controller\index{memory!controller}
    down a speed grade.
\end{itemize}

\begin{figure}[!ht]
\centering
\includegraphics[width=0.45\linewidth]{08_board_populated_hl}
\caption{G.SKILL Ripjaws S5 64\,GB DDR5-6000}\label{fig:memory}
\end{figure}


% Offset pulls the note up beside the bullets. Anchored where it sits, it
% collides with the figure caption that follows, which LaTeX does not warn about.
\marginnote[-4\baselineskip]{\newthought{Why 64 and not 128.} We will be doing fine with 64GB. 
The one real case for a large kit is
offloading a model too big for total VRAM into system memory\index{memory!offloading}, but then again you could just add another 
3090 and have 48GB of VRAM to play with.}

\newthought{Install it} by seating one module in the second slot of each channel\index{memory!channel}, follow the marks on the motherboard or its manual, 
until the clips snap up (Figure~\ref{fig:memory}).





% ==============================================================================================
% STORAGE
% ==============================================================================================

% No \newpage here: storage flows onto the page the memory subsection opened, and
% it puts two items in that margin, its datasheet citation and the capacity note.
% Those two set last on the page, so verify both still render after any edit above.
\needspace{9\baselineskip}
\subsection{Storage}

\newthought{One NVMe drive}\index{NVMe}\index{SSD|see{NVMe}}, for non volatile memory express, 2~TB in the M.2~2280\index{M.2 2280}
size and wired PCIe~4.0~x4, holding both the system and the model library
\cite{gigabyteB650EagleAX}:



\begin{itemize}[topsep=0.25em, itemsep=0.4em, parsep=0em]
    \item For speed it belongs in the first M.2 socket, the only one wired straight
    from the processor. A quantised model\index{quantization} is tens of gigabytes read start to finish
    every time it loads, and that socket is what sets how long we wait on a model swap\index{model swapping}.
    \item For width the socket is PCIe~5.0~x4 while our drive is Gen4~x4\index{PCIe!link width}, so it runs
    at its own full width, around 7.4~GB/s, which is already fast but also leaves headroom left in the socket for a
    faster drive later.
    \item We go for single storage as the second socket runs only PCIe~4.0~x2 and would halve
    this drive's lanes, and the third sits behind the chipset, sharing its uplink
    with network and USB.
\end{itemize}

\begin{figure}[!ht]
\centering
\includegraphics[width=0.45\linewidth]{12_board_complete_hl}
\caption{NVMe Gen4 x4 2\,TB, M.2 2280}\label{fig:nvme}
\end{figure}


% Offset pulls the note up beside the bullets: anchored after the itemize it sets
% level with "Fit the drive" and leaves the whole upper margin empty.
\marginnote[-6\baselineskip]{\newthought{The terabytes are a use-case call}. It leaves
room for a few dozen quantised models\index{model library} beside the system, which is plenty for pure
inference on a handful of models. If you need more models or upgrade regularly pruning will become a housekeeping
chore early one, which is best practice anyway in my opinion.
Datasets, fine tuning\index{fine-tuning} runs and their checkpoints on the other hand quickly will become an issue with 2TB, so size the
drive to what you actually intend to do.}


\newthought{Fit the drive} in the first socket, in at an angle, pressed flat, and
held down with its single screw or latch (Figure~\ref{fig:nvme}).



% ==============================================================================================
% THE FRAME
% ==============================================================================================

\needspace{9\baselineskip}
\section{The Frame}\index{open frame}

\newthought{An open GPU rack frame}. Any aluminium skeleton will do, 
a flat bed for the board and a rail above it
that holds cards upright.

\marginnote{\newthought{Cases are expensive.}\index{case} The downside is that 
the rig is open to dust and to hands, so if you can not move it out of the way - reconsider a case.
Make sure it holds your motherboard and parts in any case. Yep, pun intended.}

\begin{itemize}[topsep=0.25em, itemsep=0.4em, parsep=0em]
    \item For the card, an open bed means it hangs in room air\index{airflow} rather than behind a
    panel, so it runs cooler than the same card in a mid tower.
    \item For growth, the rail makes it easy to swap and add cards easily.
    \item Of course this also gives us a server style look, which I appreciate in this project.
\end{itemize}

\begin{figure}[!ht]
\centering
\includegraphics[width=0.77\linewidth]{14_frame_board_mounted_hl}
\caption{Open GPU rack frame}\label{fig:frame}
\end{figure}


\newthought{Mount it now.} Assemble the frame, screw the standoffs\index{standoffs} into the bed, and
lower the populated board onto them. I directly screwed this whole thing into a wall,
but you can also put it on a shelf - if you prefer to not go rogue in your flat. The frame's two pin switch is now 
wired onto \texttt{PW} on \texttt{F\_PANEL}\index{front panel header}.




% ==============================================================================================
% THE GRAPHICS CARD
% ==============================================================================================

\needspace{9\baselineskip}
\section{The Graphics Card}

\marginnote{\newthought{Why the 3090} and not:\index{RTX A6000}
\begin{itemize}[leftmargin=*, topsep=0.25em, itemsep=0.1em, parsep=0em]
    \item A6000, 48~GB: Near ten thousand euro (in 2026).
    \item 5090\index{RTX 5090}, 32~GB: More bandwidth, but no NVLink.
    \item 4090\index{RTX 4090}, 24~GB: The same VRAM at twice the price.
    \item Mac or DGX~Spark\index{DGX Spark}: Unified memory\index{unified memory} crawls on a dense 70~billion parameter
    model.
    \item 3090 Founders Edition\index{Founders Edition}: wants NVIDIA's twelve pin adapter\index{twelve pin adapter} rather than plain
    eight pin cables.
\end{itemize}}

\newthought{A RTX~3090, 24~GB} seated straight in the board's top slot:

\begin{itemize}[topsep=0.25em, itemsep=0.4em, parsep=0em]
    \item For capacity, 24~GB of GDDR6X\index{GDDR6X} VRAM on a 384~bit bus\index{memory!bus width}, which is the limiting factor for model size and
    model context\index{context window}.
    \item For token speed, the 936~GB/s of memory bandwidth\index{memory!bandwidth}\cite{nvidiaRTX3090}
    matters more than the
    10\,496 CUDA cores\index{CUDA cores} while picking a GPU.
    \item For the link, the top slot is the one PCIe~4.0~x16 connection in this
    machine, as long as you go with a single card.
\end{itemize}



\begin{figure}[!ht]
\centering
\includegraphics[width=0.75\linewidth]{15_gpu_seated_hl}
% tufte's caption takes [short][offset]; a positive offset lowers it. Sets the
% caption beside the lower part of the card instead of level with its top edge.
\caption[][12\baselineskip]{RTX~3090 24\,GB}\label{fig:gpu}
\end{figure}

\newthought{Fit the card} into the top slot, pressed down evenly until the retention clip snaps shut, fasten its bracket, and leave its power cables until the supply is
mounted. The multi-card rig extensions this motherboard is able to hold will be detailed out in a later section. This in itself is a powerful 
machine.

\marginnote{\newthought{Buying used.}\index{thermal pads}\index{used market} Buy only from a seller who
shows a test of the card under load, and budget a thermal pad replacement. For the used market, a three fan model with two eight pin connectors, an
EVGA FTW3 Ultra, MSI Suprim or Gaming, Gigabyte Gaming~OC, or ASUS~TUF.}

% Placed after the closing paragraph, not after "Fit the card": anchored there its
% caption sets on top of the Buying used marginnote, which LaTeX does not warn about.





% ==============================================================================================
% THE POWER SUPPLY
% ==============================================================================================

\needspace{9\baselineskip}
\section{The Power Supply}

\newthought{An ASRock Phantom Gaming PG-1600G}, 1600~watt against the roughly
500~watt this machine draws, bought for the headroom\index{ASRock PG-1600G}\index{PSU|see{power supply}}\index{power!budget}\cite{asrockPG1600G}:

\marginnote{\newthought{Get solar power for your idle wattage.}\index{Balkonkraftwerk}\index{solar power}\index{power!idle draw}
The floor runs day and night whether or not anyone prompts the machine, and most of it
is the platform rather than the card. A steady draw is what a balcony solar setup
covers best: panels and a capped inverter suit a constant load, not the spikes this
supply was bought to survive. Put a meter\index{power meter} on the wall socket before sizing any of it.}

\begin{itemize}[topsep=0.25em, itemsep=0.4em, parsep=0em]
    \item For the graphics card, a 3090 averages 350~watt and spikes far above it,
    which is what trips undersized supplies mid run\index{power!transient spikes}. ATX~3.1\index{ATX 3.1} rates this unit for
    220 percent of that, and the 1600~watt rating is what we need for four potential cards.
    \item For the board, one 24~pin main cable\index{ATX 24 pin} and one processor cable, 4+4~pin,
    both modular\index{modular cables}, so only what we use leaves the unit.
    \item Eight PCIe cables\index{PCIe!power cable} ship in the box. We wire two, and the rest is already
    there for later cards.
\end{itemize}

\begin{figure}[!ht]
\centering
\includegraphics[width=0.81\linewidth]{19_psu_mounted_hl}
\caption{ASRock PG-1600G, 1600\,W}\label{fig:psu}
\end{figure}


\marginnote{\newthought{What stays in the box.} The unit ships two processor cables
and this board takes one, so there is no second socket to hunt for. The two 12V-2x6\index{12V-2x6}
cables, 600~watt each, are for 40 and 50 series cards, while a 3090 takes plain eight
pin. Twelve SATA\index{SATA power} and six Molex\index{Molex} wait for spinning storage or a powered riser.}

\newthought{Wire it last.} Mount the supply on the frame, then run the cables in this
order:

\begin{enumerate}[topsep=0.25em, itemsep=0.4em, parsep=0em]
    \item The 24~pin main cable (often it is labelled MB) into \texttt{ATX}, the widest socket on the board,
    sitting on the long edge alongside the memory slots. Press until the latch clicks.
    \item The 4+4~pin processor cable, from the socket marked \texttt{CPU} on the
    supply into \texttt{ATX\_12V}\index{ATX 12V} on the board, on the short edge past the cooler,
    beside the VRM\index{VRM} heatsink nearest the rear I/O panel.
    \item The two PCIe cables to the card's own two eight pin connectors\index{eight pin connector}, along its
    top edge, one cable each and never daisy chained\index{daisy chaining}.
\end{enumerate}



\newthought{Leave the mains lead unplugged.} The rig is now fully assembled, wired, and ready to be powered on for the first time.
Skip to \hyperref[ch:engine]{the next chapter} for the firmware setup and the first
probe of the rig.









% ==============================================================================================
% CARDS TWO TO FOUR
% ==============================================================================================

\newpage
\section{Multi-Card Rig}\index{multi-card rig}

\newthought{One to three more RTX~3090s} find a place on this rig
design\cite{gigabyteB650EagleAX}. A second card
buys memory for larger models, while the third and fourth cards buy additional throughput.

\needspace{9\baselineskip}
\subsection{Card Two}\index{NVLink}

\marginnote{\newthought{What NVLink is.} A bridge that clips onto the top edge of
two cards and carries 112.5~GB/s bidirectional traffic on
the 3090.}

\newthought{Stay on the 3090.} It is the only GeForce~30 card
with NVLink, so a matched pair takes a bridge and pools\index{memory!pooling} to 48~GB for one model. 
Mixed cards never pool.

\newthought{The rest is mechanical.} Drop the card into the next x16 sized
position, wired PCIe~3.0~x1, on a plain riser\index{riser}\index{mining}. 
Riser spacing is yours to set, so match it to
the bridge you buy, get it fixed on your frame, and run its own two cables from the supply.

\needspace{9\baselineskip}
\subsection{Cards Three and Four}

\newthought{Add these with increased workload.} Two x16 sized positions remain, one
from the processor and one from the chipset. Both are wired PCIe~3.0~x1\index{PCIe!x1 lane}: full length
connector, a single lane\index{PCIe!lanes} behind it, about 1~GB/s only.
Keep it at one model per (linked) card(s). A model split across cards
parallelism\index{tensor parallelism} moves activations every token and starts to crawl. So the idea is to have the same model loaded
    on each card, and let the engine\index{engine} round robin\index{round robin} between them for throughput.

\marginnote{\newthought{On power.} Four cards at 350~watt plus the
platform reach roughly 1550~watt, which becomes uncomfortable against our supply's 1600 limit.
Undervolted\index{undervolting} near 0.875~volt each card draws about 280 and the whole rig lands near
1270, which makes it comfortable again.}

\newthought{4 Cards will definitely get you Fans.} Zip tie a bank of case
fans\index{case fans} along the rail, undervolt before the fourth cars goes in.



% ==============================================================================================
% THE MINIMAL VIABLE RIG
% ==============================================================================================

% The second thread, for a reader who wants every step of this book without the
% silicon. Same three newthought shape as a part section, same scope rule: it stops
% at the assembled, unpowered board. Everything above the metal is identical, which
% is the whole argument, and it lives in the margin rather than the main column.
\newpage
\section{The Minimal Viable Rig}\index{Raspberry Pi}\index{LPDDR4X}

\newthought{A Raspberry Pi~5 with 8~GB}, (go 16 GB if you can) an active cooler\index{active cooler}, the official
5~volt 5~amp USB-C supply\index{USB-C power} and a 64~GB microSD card\index{microSD}\cite{raspberrypiPi5}, 
will enable you to learn just as much on a tighter budget.

\marginnote{\newthought{The On-A-Budget-Thread.} Same operating system, same
container runtime, same engine binary on the same port, same clients pointed at it.
Only the model and performance shrinks.}

\newthought{The 8~GB is shared}\index{unified memory}: processor and graphics core read the same
LPDDR4X, so it is system and video memory at once. A four billion parameter model at
four bit quantisation\index{quantization} takes near 2.5~GB and leaves five for the operating system and
containers. What it costs mainly is speed as 
LPDDR4X-4267 moves roughly 17~GB/s\index{memory!bandwidth} where a 3090 moves 936, and generation reads
the whole model once per token, so answers arrive at single digit tokens per
second\index{tokens per second} rather than dozens. But hey we don't care about speed for learning, 
and buying hardware later on is the easy part.

\newthought{Build it in five minutes.} Clip the active cooler onto the board and plug
its fan into the header beside it, write the operating system image onto the microSD
card and seat the card in the slot on the underside, and leave the USB-C supply out of
the wall. A complete machine, wired and never switched on, in one hand.


% ==============================================================================================
% BRIDGE
% ==============================================================================================
\vfill
\newthought{The metal is finished}, whichever of the two machines is on
your bench. What it still lacks is a first power on and a handful of firmware values
that make a card addressable, a probe that proves every part is really on the bus,
and then an operating system, a driver, and weights loaded into memory together with some orchestration. 
Let's get on it.
